\section{Slanted Asymptotes and Other Asymptotics}
\label{sec:slanted-asymptotes}

Section~\ref{sec:asymptotes} introduced horizontal and vertical asymptotes.
What if a graph keeps rising but gets closer and closer to a tilted line,
or even to a curve? The same idea still works: use a limit to measure the
error in the approximation.

\subsection{Slanted Asymptotes}

The line $y=mx+b$, with $m\ne0$, is a \textbf{slanted asymptote}
(also called an \textbf{oblique asymptote}) as $x\to\infty$ if
$$\lim_{x\to\infty}\bigl(f(x)-(mx+b)\bigr)=0.$$
In other words, the vertical gap between the graph and the line approaches
0. Use $x\to-\infty$ to describe the left end of the graph instead.

To find the line, compute the slope first and then the intercept:
$$m=\lim_{x\to\infty}\frac{f(x)}{x},\qquad
b=\lim_{x\to\infty}\bigl(f(x)-mx\bigr).$$
Both limits must be finite, and the value of $m$ from the first limit
goes into the second. If $m=0$, this is just a horizontal asymptote.
Check $x\to-\infty$ separately; it need not give the same line.

For rational functions, polynomial division often does the work for us.
In general, a rational function has a slanted asymptote when the
numerator's degree is exactly one more than the denominator's: divide,
and keep the linear quotient.
\begin{exampleblock}
Polynomial division gives
$$f(x)=\frac{x^2+1}{x-1}=x+1+\frac{2}{x-1}.$$
Since
$$\lim_{x\to\pm\infty}\bigl(f(x)-(x+1)\bigr)
=\lim_{x\to\pm\infty}\frac{2}{x-1}=0,$$
the slanted asymptote is $y=x+1$ in both directions.
\end{exampleblock}

\begin{exampleblock}
Finding a slope is not enough. If $f(x)=x+\sqrt{x}$, then
$$\lim_{x\to\infty}\frac{f(x)}x=1,\qquad
\lim_{x\to\infty}\bigl(f(x)-x\bigr)=\infty.$$
So there is no finite intercept $b$, and no slanted asymptote.
\end{exampleblock}

\subsection{Other Asymptotics}

We can also look for a simple curve $y=g(x)$ such that
$$\lim_{x\to\infty}\bigl(f(x)-g(x)\bigr)=0.$$
\begin{exampleblock}
Polynomial division gives
$$\frac{x^3+1}{x-1}=x^2+x+1+\frac{2}{x-1}.$$
The last term approaches 0 as $x\to\pm\infty$, so the graph approaches
the parabola $y=x^2+x+1$. This is sometimes called a
\textbf{polynomial asymptote}.
\end{exampleblock}

Sometimes we only want to compare a function's leading growth. We write
$$f(x)\sim g(x)\quad\text{as }x\to\infty
\qquad\text{if}\qquad
\lim_{x\to\infty}\frac{f(x)}{g(x)}=1,$$
where $g(x)\ne0$ for all sufficiently large $x$. This says that the
\textit{relative} error goes to 0, not necessarily the vertical gap.
\begin{exampleblock}
$x+\sqrt{x}\sim x$, even though their difference grows without bound.
\end{exampleblock}
Keep the two tests straight: use a difference to check
whether the graphs get close, and a ratio to compare their leading growth.
